\section{What the pieces are, and what can realize them}\label{sec:realizations}

\subsection{The two proposals, stated precisely}\label{sec:proposals}

Before naming the pieces we record what the two proposals of the introduction
actually require, since the inherited note \cite{LoewnerNote} found that both
had been stated too strongly.

\paragraph{What is abelian.} By Lemma~\ref{lem:multiplicative} the map
$H\mapsto P_L(h*\cdot)P_L$ is a homomorphism from the commutative algebra of
causal symbols onto $\mathcal A_L$, so every $V_{\omega,L}$ and every
$G_{\omega,L}$ lie in one commutative algebra, for all $\omega$ at fixed $L$.
What is abelian is \emph{translation invariance along the line}: each shift
acts on $I_L$ by a kernel depending only on the separation $x-y$, with the
prime atoms at fixed separations, and two functions of the line momentum
commute. Read as a condition on a realization, the insertions implementing the
$\omega$-deformation must act as functions of a single operator, uniformly
along the line. That is a selection of a Cartan-like direction, not an
impossibility.

\paragraph{The critical path is ordered.} The endpoint generator
$B_L=\frac12\delta_{L/2}\otimes\mathrm{ev}_{L/2}$ of the length direction
satisfies $G_{\omega,L}B_L=0\neq B_LG_{\omega,L}$ \cite[Proposition~9.3]{WilsonLines}.
Along a path $t\mapsto(\omega(t),L(t))$ the transport
$\mathcal P\exp\int(\dot LB_L-\dot\omega G_{\omega,L})\dd t$ is a genuinely
ordered product; only its value is path-independent, because the connection is
flat.

\paragraph{Flatness is a tautology.} If $V(\alpha)$ is any smooth family with
$dV=\mathcal AV$ and dense range, then $0=d^2V=(d\mathcal A-\mathcal A\wedge\mathcal A)V$
forces the curvature to vanish. In the non-Abelian Stokes variation of an open
line $W[C]$, the field strength enters as the \emph{coefficient} $\mathcal A$
of the bulk deformation, and $\mathcal A$'s curvature on deformation space is
zero by the same argument. So flatness constrains no realization; the object
to match is the coefficient $G_{\omega,L}$ itself, together with the one-sided
commutation with the endpoint term. What remains as the obstruction is the
dichotomy: any realization must be exactly critical.

\paragraph{The Loewner affinities.} Whole-plane Loewner evolution grows a hull
$K_t$ from $0$ with logarithmic capacity $t$, so Loewner time is
log-conformal-radius, the arithmetic coordinate $x=\log r$; the stabilizer of
the marked points $0,\infty$ is the group of dilations and inversions, the
exact symmetry group of the Weil form; the exterior Loewner equation
$\partial_tg_t(z)=g_t(z)\frac{e^{iU_t}+g_t(z)}{e^{iU_t}-g_t(z)}$ is a tip
equation with one scalar coefficient, like the endpoint variation
$\partial_LV_{\omega,L}=B_LV_{\omega,L}$ \cite[Proposition~9.1]{WilsonLines};
the chain is a path-ordered flow of Witt-algebra vector fields
\cite{BauerBernard,Lawler}; and composition with a univalent self-map is a contraction
on Hardy-type spaces by Littlewood subordination \cite{Shapiro}. The proposal
is therefore well motivated, and the question it poses is concrete: can such a
chain produce the kernel $k_\omega$, or at least the positive part $\kt_\omega$?

\subsection{The archimedean factor is Bessel additivity}

\begin{proposition}[Bessel additivity]\label{prop:bessel}
Let $Z_\delta$ denote independent $\mathrm{Gamma}(\frac\delta2)$ variables ---
the value at any fixed time, up to the common factor $2t$, of a squared Bessel
process of dimension $\delta$ started at $0$. Then, with $U=Z_a/(Z_a+Z_{2\omega})$,
\begin{equation}\label{eq:bessel}
 U\sim\mathrm{Beta}\Big(\frac a2,\omega\Big),\qquad
 K^\Gamma_\omega(p)=\pi^\omega\frac{\Gamma(\frac a2)}{\Gamma(\frac b2)}\E\big[U^{p/2}\big],
 \qquad Z_a+Z_{2\omega}\stackrel d=Z_b .
\end{equation}
The archimedean delay is $x=-\frac12\log U$, i.e.\ $(r_1/r_2)^2=U$ in the ray
variable $r=e^x$.
\end{proposition}
\begin{proof}
$G_\alpha/(G_\alpha+G_\beta)\sim\mathrm{Beta}(\alpha,\beta)$ and
$G_\alpha+G_\beta\sim G_{\alpha+\beta}$ for independent Gamma variables are
classical; $\frac a2+\omega=\frac b2$. The additivity of squared Bessel
processes in the dimension is Shiga--Watanabe \cite{ShigaWatanabe,RevuzYor}.
The moment identity is \eqref{eq:betamoments}.
\end{proof}

Three remarks. The two Bessel dimensions $a,b$ are the two exponents $e^{ax}$,
$e^{-bx}$ of the rational factor: one set of numbers, $\frac12\pm\omega$,
because both come from $s\pm\omega$. The dimensions are fractional and below
one, so there is no sphere and no angle; the Beta variable is the additivity
fraction, and the archimedean ``quarter'' of the delay test
\cite[Section~11]{WilsonLines} is $\frac a2$ at $\omega=0$, half a Bessel
dimension. And the statement is about Bessel processes, which drive radial
Loewner chains but are not chains: whether a conformal-radius or hitting law
of a radial chain driven by a Bessel process of dimension $b$ has $U$ as a
squared-radius ratio is open (Section~\ref{sec:outlook}).

\subsection{A Loewner chain cannot produce the comb}\label{sec:nocomb}

There are two ways to identify the arithmetic coordinate with a Loewner
variable, and neither produces the comb.

\begin{lemma}[Rotation averages of composition operators are dilations]\label{lem:rotation}
Let $\phi$ be analytic near $0$ with $\phi(0)=0$ and let $h$ be analytic. Then
\begin{equation}\label{eq:rotation}
 \frac1{2\pi}\int_0^{2\pi}h\big(e^{i\theta}\phi(e^{-i\theta}z)\big)\dd\theta=h\big(\phi'(0)z\big).
\end{equation}
Consequently, for a random conformal map $g$ with rotation-invariant law and
deterministic $g'(0)$, $\E[h\circ g]=h(g'(0)\,\cdot)$ for every analytic $h$;
the same holds at $\infty$ for exterior maps $g(z)=e^{-t}z+O(1)$.
\end{lemma}
\begin{proof}
It suffices to take $h=z^n$. Then $e^{in\theta}\phi(e^{-i\theta}z)^n=e^{in\theta}\sum_{k\geq n}b_ke^{-ik\theta}z^k$
and the average keeps only $k=n$, with coefficient $b_n=\phi'(0)^n$.
\end{proof}

So if the arithmetic coordinate is Loewner time, the expected composition
operator of a rotation-invariant chain with the capacity normalization acts on
$H^2$, on the Dirichlet space, on any space of analytic functions, as
$h\mapsto h(e^{-t}\cdot)$: a pure delay by $t$, symbol $e^{-tp}$, carrying no
information about the chain. The shape of the curve lives in non-analytic
observables --- $|g|$, harmonic measure, $\E|g'|^q$ --- and a Loewner
realization of the transfer would have to act on those. The natural one is the
radial process $\log|g_t(z)|$, and for it the other identification applies.

\begin{proposition}[Far-field transport is a pure delay]\label{prop:farfield}
Let $K_t$ be a whole-plane Loewner hull of logarithmic capacity $t$,
deterministic or random, with any number of tips, and
$g_t:\hat\C\setminus K_t\to\hat\C\setminus\overline{\mathbb D}$ its normalized
map, $g_t(z)=e^{-t}z+O(1)$. Then $\log|g_t(z)|=\log|z|-t+O(|z|^{-1})$ uniformly
in $\arg z$. Hence if the induced transition kernel $P_t(y,\dd y')$ from
$y=\log|z|$ to $y'=\log|g_t(z)|$, averaged over $\arg z$ and over the chain's
law, is translation-invariant in $y$, it equals $\delta_{y'=y-t}$. In
particular no Loewner chain induces a translation-invariant radial kernel with
atoms at $y'=y-\log n$, $n\geq2$.
\end{proposition}
\begin{proof}
The normalization at infinity is the definition of the capacity
parametrization; translation invariance forces $P_t(y,\cdot)$ to equal its
$y\to\infty$ limit.
\end{proof}

Away from the far field the kernel depends on the distance to the hull and is
not translation-invariant; and in either case $t\mapsto g_t(z)$ is continuous
even for a c\`adl\`ag driving function, since the Loewner equation is an ODE
whose field is smooth in $g$ away from the circle, so a jump of the driving
process moves the tip, not the transported points. Echoes at the fixed
separations $\log n$ with position-independent weights are not a Loewner
phenomenon in any reading. What Loewner or Bessel processes can supply is the
archimedean factor alone.

\subsection{What the comb is}\label{sec:hecke}

On functions of the arithmetic coordinate let $\mu_nf(x)=f(x-\log n)$. The
$\mu_n$ are isometries of $L^2(\R,\dd x)$ with $\mu_n\mu_m=\mu_{nm}$, and the
comb operator is
\begin{equation}\label{eq:heckeop}
 C_\omega=\sum_{n\geq1}\ct_n\,\mu_n,\qquad
 \ct_n=n^{-1/2}\big(n^{\omega}*\mu(n)n^{-\omega}\big)=n^{\omega-\frac12}\prod_{p\mid n}(1-p^{-2\omega}),
\end{equation}
a Dirichlet series in the multiplicative semigroup $\N^\times$ acting on the
ray by dilations: the image of $\zeta(s-\omega)/\zeta(s+\omega)$ under the
representation $n^{-s}\mapsto\mu_n$, with the $n^{-1/2}$ making it unitary on
$L^2(\dd x)$. This is the Hecke algebra of $\N^\times$ in the Bost--Connes
system \cite{BostConnes}, which \cite[Section~10.4]{WilsonLines} had already
named as the structure ``invariant under $r\mapsto\lambda r$ but distinguishing
the ratios $p^m$'' that the prime atoms require; the decomposition says exactly
which element of that algebra the transfer contains, with which coefficients,
at each shift. Its local factor at $p$, $\frac{1-p^{-\omega}z_p}{1-p^{\omega}z_p}$
with $z_p=e^{-s\log p}$, is a one-delay feedback loop of length $\log p$ with
loop gain $p^{\omega}$: its kernel before the $n^{-1/2}$ is
$\delta_0+(1-p^{-2\omega})\sum_{k\geq1}p^{k\omega}\delta_{k\log p}$, and a
traversal of the loop multiplies by the scalar $p^\omega$, which is the abelian
monodromy at the puncture of \cite[Section~9.4]{WilsonLines}, made explicit.
To first order in $\omega$ the coefficients are $2\omega\Lambda(n)n^{-1/2}$,
the von Mangoldt function, and at second order the composite weights appear:
$\ct_6=\omega^2\cdot4\log2\log3/\sqrt6+O(\omega^3)$.

\subsection{The endpoint \texorpdfstring{$\omega=\frac12$}{omega = 1/2} is the modular surface}\label{sec:endpoint}

At $\omega=\frac12$: $a=0$, $b=1$, $\frac{p+a}{p-a}=1$, so
\begin{equation}\label{eq:endpoint}
 \Kh_{1/2}(p)=\Kt_{1/2}(p)=\frac{\Lambda(p)}{\Lambda(p+1)}=\varphi\Big(\frac{p+1}2\Big),
 \qquad
 \varphi(s')=\sqrt\pi\,\frac{\Gamma(s'-\frac12)\,\zeta(2s'-1)}{\Gamma(s')\,\zeta(2s')}
 =\frac{\Lambda(2s'-1)}{\Lambda(2s')},
\end{equation}
where $\varphi$ is the scattering matrix of the Eisenstein series of
$PSL(2,\mathbb Z)\backslash\mathbb H$ \cite{Iwaniec}, the coefficient of
$y^{1-s'}$ in the constant term of $E(z,s')$; its pole at $s'=1$ is $p=b=1$.
The comb weights are $\ct_n=\prod_{p\mid n}(1-p^{-1})=\varphi_{\mathrm{Euler}}(n)/n$,
the density of reduced fractions with denominator $n$ --- the count of cosets
in the cusp's Fourier expansion that produces $\zeta(2s'-1)/\zeta(2s')$ in the
first place. The Beta part degenerates: $\mathrm{Beta}(0,\frac12)$, and
$k^\Gamma_{1/2}(x)=2(1-e^{-2x})^{-1/2}$ does not decay, which is the pole of
$K^\Gamma_{1/2}$ at $p=0$. And
\begin{equation}\label{eq:laxphillips}
 K_{1/2}=B_1\,\varphi\Big(\frac{p+1}2\Big)=\frac{p-1}{p+1}\,\varphi\Big(\frac{p+1}2\Big):
\end{equation}
the transfer at the endpoint is the Eisenstein scattering matrix with its pole
at $s'=1$ removed by one Blaschke factor, which is what Lax and Phillips do
when they pass to the orthocomplement of the constant function
\cite{LaxPhillips} (a reading; the identity is exact). So the Markov part has,
at the endpoint, a classical geometric realization in which the comb is a
coset sum and the correction is the removal of the residual spectrum. For
$\omega<\frac12$ Eisenstein constant terms give only unit offsets in $2s'$, so
the interior of the family is not realized by any congruence-subgroup
scattering matrix; this was already observed in \cite{WilsonReview}. All of
\eqref{eq:endpoint}--\eqref{eq:laxphillips} and the totient identity are
checked numerically in Appendix~\ref{app:numerics}.

\subsection{Summary of the dictionary}

\begin{center}
\small
\begin{tabular}{@{}lll@{}}
\toprule
Piece & Symbol & Realization \\
\midrule
Archimedean delay & $K^\Gamma_\omega=\pi^\omega\Gamma(\frac{s-\omega}2)/\Gamma(\frac{s+\omega}2)$ & Bessel additivity in dimensions $a,2\omega\to b$ (Prop.~\ref{prop:bessel}) \\
Comb & $K^\zeta_\omega=\zeta(s-\omega)/\zeta(s+\omega)$ & Hecke operator $\sum\ct_n\mu_n$ of $\N^\times$; not Loewner (Prop.~\ref{prop:farfield}) \\
Correction & $R_\omega=B_b/B_a$; cleanly $B_b$ & Forced by unimodularity (Prop.~\ref{prop:forced}); no arithmetic \\
Endpoint $\omega=\frac12$ & $K_{1/2}=B_1\varphi$ & Modular surface; Lax--Phillips (\S\ref{sec:endpoint}) \\
The whole & $K_\omega=B_b\Kh_\omega$ & Positive measure minus twice its EMA (Prop.~\ref{prop:blaschke}) \\
\bottomrule
\end{tabular}
\end{center}
